% !TEX root = Thai Red Chicken Meatballs Meal Prep.tex
\documentclass[14pt]{extarticle}
\usepackage[utf8]{inputenc}
% --- REMOVE PAGE NUMBERING ---
\pagestyle{empty}
% --- MARGIN SETUP ---
\usepackage[left=0.75in, right=0.75in, top=0.5in, bottom=0.5in]{geometry}
% --- LIST SETTINGS ---
\usepackage{enumitem}
\setlist[itemize]{leftmargin=1em}
% Labels are centered in a fixed-width box so plain bullets line up with the
% wider prep icons (\veg, \measure) below.
\SetLabelAlign{center}{\makebox[\labelwidth][c]{#1}}
\setlist[itemize]{leftmargin=1em, itemsep=0.2em, label={\raisebox{-0.1ex}{\large\textbullet}},
    align=center, labelwidth=0.9em, labelsep=0.4em}
\setlist[enumerate]{leftmargin=1.5em, label=\textbf{\arabic*.}, itemsep=0.25em}
\usepackage{parskip}
% --- CONTROL SECTION FONT SIZE ---
\usepackage{titlesec}
\titleformat{\section}{\large\bfseries}{}{0em}{}
\titleformat{\subsection}{\normalsize\bfseries}{}{0em}{}
\titlespacing*{\subsection}{0pt}{0.75em}{0.25em}
\titlespacing*{\section}{0pt}{1em}{0.5em}
% --- FONT SETUP ---
\usepackage{helvet}
\renewcommand{\familydefault}{\sfdefault}
% --- CUSTOM COMMAND FOR TITLE ---
\newcommand{\makeRecipeTitle}[1]{
    {\noindent \Large \textbf{#1}}
    \par \vspace{-0.75em}
}
% --- CUSTOM COMMAND FOR SERVINGS ---
% Usage: \recipeServings{4--5} or \recipeServings{6 (one burrito each)}
% Place directly after the first \makeRecipeTitle on page 1.
\newcommand{\recipeServings}[1]{
    {\noindent \itshape Serves #1 \par}
    \vspace{-0.75em}
}
% --- PREP ICONS ---
% Usage: \item[\veg] Red Bell Pepper – 1, sliced
%        \item[\measure] Fish Sauce – 1 1/2 tbsp (22ml)
% Replace the bullet on lines to handle before cooking starts: produce to cut
% (carrot), and spices, seasonings, sauces, and pastes with a measured amount
% (measuring spoon, drawn with TikZ). "To taste" lines get neither.
\usepackage{fontawesome5}
\usepackage{tikz}
\newcommand{\veg}{\raisebox{0.25ex}{\small\faCarrot}}
\newcommand{\measure}{\tikz[x=0.9em, y=0.9em,
        baseline={([yshift=-0.36em]current bounding box.center)}]{
    \begin{scope}[rotate=45]
        % Bowl with a rim highlight, and a flat handle with a hanging hole.
        \fill (0,0) ellipse [x radius=0.33, y radius=0.29];
        \draw[white, line width=0.035em, line cap=round]
            (-0.1,0.173) arc[start angle=120, end angle=210, radius=0.2];
        \fill (0.2,-0.07) -- (0.86,-0.105) arc[start angle=-90, end angle=90, radius=0.105]
            -- (0.2,0.07) -- cycle;
        \fill[white] (0.86,0) circle (0.04);
    \end{scope}
}}
% --- CUSTOM COMMAND FOR SOURCE CREDIT ---
% Usage: \recipeSource{Creator -- Video Title}{https://...}
% Place just before \end{document}; pins the credit to the bottom of the last page.
\usepackage[hidelinks]{hyperref}
\urlstyle{same}
\newcommand{\recipeSource}[2]{
    \vfill
    {\noindent \footnotesize \textbf{Source:} #1 \\ \url{#2}}
}


\begin{document}

% ================= PAGE 1 =================
\makeRecipeTitle{Thai Red Chicken Meatballs Meal Prep}
\recipeServings{5 (meal-prep containers)}

\section*{Ingredients}

\textbf{The Meatballs}
\begin{itemize}[itemsep=0.1em, label={\raisebox{-0.1ex}{\large\textbullet}}]
      \item Ground Chicken – 1.5 lbs (700g)
      \item[\cut] Zucchini – 6.7 oz (190g), grated
      \item[\cut] Red Chili – 1, diced
      \item[\cut] Ginger – 2 tbsp (20g), grated
      \item[\cut] Garlic Cloves – 3, grated
      \item[\cut] Cilantro – 1/2 cup (10g), chopped
      \item[\cut] Green Onions – 3, sliced
      \item Egg – 1 large
      \item[\measure] Soy Sauce – 2 1/2 tbsp (40ml)
      \item[\cut] Lime – 1, zested (save it for juicing)
      \item Salt and Pepper – to taste
\end{itemize}

\textbf{The Rice}
\begin{itemize}[itemsep=0.1em, label={\raisebox{-0.1ex}{\large\textbullet}}]
      \item Jasmine or Long Grain White Rice – 1 1/2 cups (300g), rinsed
      \item[\measure] Cold Water – 2 1/3 cups (550ml) for jasmine, or 3 cups (710ml) for long grain
      \item Salt – to taste
\end{itemize}

\newpage
\makeRecipeTitle{Thai Red Chicken Meatballs Meal Prep}
\section*{Ingredients (Continued)}

\textbf{The Sauce}
\begin{itemize}[itemsep=0.1em, label={\raisebox{-0.1ex}{\large\textbullet}}]
      \item[\measure] Thai Red Curry Paste – 3 tbsp (45g)
      \item[\cut] Green Onions – 3, sliced
      \item[\cut] Ginger – 2 tbsp (20g), grated
      \item[\cut] Garlic Cloves – 3, grated
      \item Coconut Milk – 1 can, 14 oz (400ml)
      \item[\cut] Red Bell Pepper – 1 (190g), sliced
      \item[\measure] Chicken Stock – 1 cup (250ml)
      \item[\measure] Fish Sauce – 1 1/2 tbsp (22ml)
      \item[\measure] Soy Sauce – 1 1/2 tbsp (22ml)
      \item[\cut] Broccolini – 8.8 oz (250g), trimmed
      \item[\measure] Lime Juice – from the zested lime
      \item[\cut] Cilantro – to garnish
      \item Salt and Pepper – to taste
\end{itemize}

\section*{Notes}

\begin{itemize}
      \item \textbf{Storage:} Keeps for up to 4 days in the refrigerator, or up to 3 months in the freezer. Thaw overnight in the refrigerator, then reheat gently in a pan over medium-low heat, stirring well, or microwave for 2--3 minutes, stirring halfway, until steaming hot throughout.
\end{itemize}

\newpage
\makeRecipeTitle{Thai Red Chicken Meatballs Meal Prep}
\section*{Instructions}

\begin{enumerate}

      \item \textbf{Prep Everything} \\
            Grate the zucchini and all the ginger and garlic, and slice all the green onions. Dice the chili, chop the cilantro, slice the bell pepper, trim the broccolini, and zest and juice the lime. Rinse the rice, and measure out everything marked with a spoon.

      \item \textbf{Make the Meatball Mix} \\
            Squeeze the excess moisture out of the grated zucchini with a paper towel; skipping this makes the mix too wet to hold its shape. Combine with the rest of the meatball ingredients and mix gently; don't overwork it.

      \item \textbf{Brown the Meatballs} \\
            Heat a large pan over medium heat with a little oil. Shape the mix into golf ball-sized meatballs and brown them in batches on all sides. They don't need to be cooked through yet. Set aside.

      \item \textbf{Cook the Rice} \\
            Add the rinsed rice to a pot with the water and a pinch of salt and bring to a boil. Cover, reduce to low, and cook for 12 minutes (15--18 minutes for long grain). Take the pot off the heat and let it rest with the lid on for 4 minutes.

      \item \textbf{Bloom the Curry Paste} \\
            In the same pan, add a little oil if needed. Fry the curry paste, green onions, ginger, and garlic until fragrant, scraping up the browned bits (1--2 minutes). Frying the paste in oil first is what blooms the spices.

      \item \textbf{Simmer the Sauce} \\
            Add the coconut milk, bell pepper, chicken stock, fish sauce, and soy sauce. Bring to a gentle simmer and let it reduce (5 minutes).

      \item \textbf{Finish the Curry} \\
            Return the meatballs to the pan and simmer (6 minutes). Add the broccolini and simmer until it's just tender but still has bite, and the meatballs reach 165$^\circ$F (4 minutes); any longer and it turns mushy over the week. Season with salt and pepper, then finish with the lime juice and extra cilantro.

      \item \textbf{Portion and Store} \\
            Divide the rice between 5 meal-prep containers (750ml each), then top with the meatballs, sauce, and broccolini. Let cool before sealing.

\end{enumerate}

\recipeSource{Chef Jack Ovens -- Thai Red Chicken Meatballs Meal Prep}{https://chefjackovens.com/thai-red-chicken-meatballs-meal-prep/}

\end{document}
